\begin{abstract}
Speculative decoding splits naturally between the edge and the cloud: a small draft model on the device proposes tokens and a large cloud model verifies them. Every verification, however, crosses the network, and in synchronous speculative decoding the device sits idle through every round trip. We present \cape{}, an edge-only controller that turns this idle time into progress. In every round, \cape{} decides how many tokens to draft on the critical path and how many to prefetch while the cloud verifies. It models decoding as a two-state Markov chain over prefetch hits and misses and prices the measured round-trip distribution, the interruption of prefetches that a rejection has made useless, and the device energy of every speculative token. It needs no offline calibration, no draft probabilities and no change to the cloud server. On a Jetson Orin Nano drafting over campus Wi-Fi for four draft/target pairs from the Llama, Qwen and Gemma families (\E{n:main} paired requests, nine methods), \cape{} is \E{x:S:sync:rng} faster than synchronous speculative decoding in all twelve pair--delay settings. When the delay changes over time, it is \E{x:S:sync:dyn:rng} faster and the fastest of all nine methods on every pair. One configuration serves every pair and delay: \cape{} is never more than \E{x:S:regret} behind the fastest method of any setting, whereas every baseline falls at least \E{x:base:regret} behind in some setting. Its energy setting matches the speed of the best fixed window chosen in hindsight under time-varying delay and uses up to \E{gemma3:dyn:E:best:svm} less device energy per token.
\end{abstract}

\begin{IEEEkeywords}
speculative decoding, edge--cloud collaboration, large language models, energy efficiency, embedded GPU
\end{IEEEkeywords}
